\section{Introduction}
\label{sec:intro}

A pure functional language promises that a value, once made, never changes. Compiled to a
machine with mutable memory, the promise has a price: an update of a large structure must either
copy it or use a persistent representation that shares the unchanged parts. The price is
avoidable when nobody can tell the difference---when no part of the program will ever look at the
old version again. The long line of work on linear and uniqueness types
\citep{wadler1990linear,wadler1991use,smetsers1994guaranteeing,barendsen1996uniqueness,devries2008simplified},
on destructive-update analysis \citep{wand2001set,cann1990sisal,aspinall2008usage}, and more
recently on reference counting with reuse \citep{ullrich2019beans,reinking2021perceus,lorenzen2023fp2}
and on modes in OCaml \citep{lorenzen2024oxidizing} all ask, in one form or another, when an
in-place write is unobservable.

This paper asks the question for one language and one data type, with an unusually strict
answer in mind. beni is a strict, pure, Elm-like language that compiles to JavaScript and is aimed
at the browser. Its single sequence type, \code{List}, is backed by JavaScript arrays. We want
every \code{List} to be a \emph{plain} JavaScript array that is written \emph{in place} by every
update the program performs, and we want the compiler to guarantee that this is safe. Concretely:

\begin{quote}
\emph{Editing a list while anything else still holds the old version is a compile error.}
Ownership is inferred, user code carries no annotations, and an accepted program writes in place,
always---in development and release builds alike---with no reference count, no run-time test and
no silent fallback to a copy.
\end{quote}

Two things make this goal different from most prior work. First, the result is an \emph{error},
not an optimisation. Run-time schemes such as Perceus \citep{reinking2021perceus}, Lean~4's
borrow inference \citep{ullrich2019beans} and Roc's in-place mutation of uniquely referenced
values \citep{teeuwissen2023roc}
update in place when a reference count says the value is unique and copy otherwise; the program is
always accepted, and its cost depends on facts the programmer cannot see. Static schemes that
\emph{optimise}---Sisal's copy elimination \citep{cann1990sisal}, Wand and Clinger's set
constraints \citep{wand2001set}---also fall back to a copy. We want the opposite contract: a
program that would need a hidden copy is rejected, with a message that names the holder of the old
version and offers the explicit copy. Second, the language has no references, no lifetimes and no
ownership vocabulary in its surface syntax, so everything must be inferred; the type-based
uniqueness systems of Clean \citep{barendsen1996uniqueness,plasmeijer2011clean} and Futhark
\citep{henriksen2017futhark} express uniqueness as attributes in types; in Futhark the programmer writes
them on every function that consumes an argument.

\paragraph{Why it matters in a JavaScript target.}
A persistent sequence in JavaScript is not free. beni's current representation (\S\ref{sec:bg-lists})
is an array that switches to a 32-way radix trie when shared writes grow large; it costs runtime
code that every program that edits a list must ship, a representation test on every read site that
can meet a trie, and a conversion whenever a value crosses into JavaScript. In a prototype study
(\S\ref{sec:bg-measure}) in which we applied by hand the transformations a static uniqueness pass
would make to compiled beni code, building an array in a fold by in-place writes ran at 0.14--0.47
times the time of the persistent representation, the steady-state ticks of a million-cell grid at
0.44--0.48 times, and the runtime support for lists shrank from 1\,514 to 706 bytes after compression
once the trie was gone. The same study found the obstacle: a UI runtime that keeps the lists it
rendered makes every list in the model shared, so no static rule could write them in place until
the runtime hands the model over.

\paragraph{Why beni can infer it.}
Three properties of the language do most of the work.
(1)~\emph{Strict, left-to-right evaluation in source order.} The order in which a read and a write
happen is the order in which they are written, so ``before the update'' is a syntactic notion, as
in Futhark and in Smetsers et al.'s classification of argument positions
\citep{smetsers1994guaranteeing}.
(2)~\emph{No automatic currying.} Every call in beni is saturated; partial application is written
as an explicit lambda. The case every uniqueness system has found hardest---a partial application
that silently captures a unique argument
\citep{barendsen1996uniqueness,devries2008simplified,lorenzen2024oxidizing}---does not arise: a
closure is a lambda written in the source, and its captures are visible to the checker.
(3)~\emph{An existing inference of per-function facts.} beni's type checker already infers two
effect bits for every function type, solves them as a directed constraint graph after each module,
and publishes them in the module interface. Ownership facts can ride in the same places.

\paragraph{Contributions.}
This paper is a design proposal under evaluation. Its contributions are:
\begin{itemize}
\item A precise statement of the property to check: \emph{uniqueness at the update site, decided by
liveness} (\S\ref{sec:model-unique}). It is strictly more permissive than Clean's occurrence-based
uniqueness: reads before a write are never sharing, and a dead alias is harmless, in the spirit of
Boyland's alias burying \citep{boyland2001alias} and Rust's two-phase borrows \citep{rfc2025}.
\item An inference algorithm (\S\ref{sec:inference}): a directional, flow-sensitive,
field-sensitive cell analysis, a per-path liveness pass, four side conditions at each update
(\Known, \Unescaped, \Dead, \Disjoint), function summaries polymorphic over the ownership
behaviour of function-typed parameters, an optimistic fixpoint for recursive groups, and
publication of summaries in module interfaces. It is linear in the size of a function body per
round, with a round count bounded by the height of a finite lattice.
\item A soundness argument (\S\ref{sec:soundness}): under ten stated assumptions about the
runtime, the core library and the platforms, an accepted program produces the same observable
trace under value semantics and under in-place semantics. We give the theorem, the key lemmas and
proof sketches; the sketch has not been independently reviewed, and we list the places we
consider weakest.
\item A treatment of the hard cases (\S\ref{sec:hard}), including closures called many times,
nested containers (with an \emph{exclusive-contents} flag for the case where two slots of a
container may hold the same value), fibers and commands, the foreign-function boundary, and a
hand-over protocol between a TEA (Elm Architecture) runtime and the program's \code{update}
function.
\item An error-message design (\S\ref{sec:errors}) that names three program points, writes the
fix, and tells \emph{real sharing} from \emph{a limit of the checker} by a derivation tag carried
on every fact, together with a stability rule for the accepted set across compiler releases.
\item An honest account of the trade-off (\S\ref{sec:discussion})---the rule rejects correct
programs to guarantee performance---and an evaluation plan (\S\ref{sec:evaluation}) that measures
the rejection rate on real code in a report-only mode before the rule is adopted. \emph{No results
of that evaluation exist yet}; we report none.
\end{itemize}

Section~\ref{sec:background} introduces beni and the representation question.
Sections~\ref{sec:model}--\ref{sec:errors} present the design. Section~\ref{sec:compiler}
describes its integration with the rest of the compiler, \S\ref{sec:discussion} its costs and
limits, \S\ref{sec:evaluation} the evaluation plan, and \S\ref{sec:related} related work.
Appendix~\ref{app:examples} works four programs through the rules.
